% =============================================================================
%  postextual/apendices.tex
%  APÊNDICES — documentos criados pelo PRÓPRIO AUTOR
%
%  Exemplos: questionários, roteiros de entrevista, manuais do usuário,
%            scripts de instalação, detalhamento de casos de teste...
%
%  Numeração: Apêndice A, Apêndice B, Apêndice C...
%  O abntex2 cuida da numeração automaticamente.
% =============================================================================

\begin{apendicesenv}

% \partapendices gera a página divisória com o título "APÊNDICES" em destaque.
% Exigida pelo abntex2 e pelas normas ABNT.
\partapendices


% -----------------------------------------------------------------------------
%  APÊNDICE A
%  Substitua o título e conteúdo pelo seu apêndice real.
%  Exemplo de uso: questionário de pesquisa com usuários.
% -----------------------------------------------------------------------------
\chapter{TODO: Título do Apêndice A}
\label{ap:questionario}

TODO: Descreva brevemente o contexto deste apêndice — por que ele
foi elaborado e como foi utilizado no projeto.

% Exemplo de questionário:
\begin{enumerate}
  \item TODO: Pergunta 1 — ex.: Com que frequência você realiza a
        atividade X?
        \begin{enumerate}[label=(\alph*)]
          \item Nunca
          \item Raramente (menos de 1 vez por semana)
          \item Às vezes (1 a 3 vezes por semana)
          \item Frequentemente (mais de 3 vezes por semana)
        \end{enumerate}

  \item TODO: Pergunta 2 — ex.: Você já utilizou alguma ferramenta
        para Y? Qual?

        \noindent\rule{\textwidth}{0.4pt} % Linha para resposta aberta

  \item TODO: Adicione as demais perguntas do seu instrumento.
\end{enumerate}


% -----------------------------------------------------------------------------
%  APÊNDICE B (descomente e adapte se necessário)
% -----------------------------------------------------------------------------
% \chapter{TODO: Título do Apêndice B}
% \label{ap:manual}
%
% TODO: Conteúdo do Apêndice B.


\end{apendicesenv}
